\documentclass{article}
\usepackage[UTF8]{ctex}
\usepackage{amsmath,amssymb,amsfonts}
\usepackage{enumitem}
\usepackage{geometry}
\geometry{a4paper, margin=1in}
% === 错题本排版优化（不参与题目标识） ===
\setlength{\parindent}{0pt}
\setlength{\parskip}{0.35em plus 0.10em minus 0.08em}
\setlength{\emergencystretch}{3em}
\setlist{itemsep=0.25em, topsep=0.40em, parsep=0pt}
\allowdisplaybreaks[2]
% === 错题本排版优化结束 ===
\begin{document}
\section*{西南交通大学2026年数学分析真题}
\section*{172．西南交通大学2026年数学分析试卷（回忆版）}

\subsection{题目1}
1．（10 分）已知 \(\left\{x_{n}\right\}\) 是无界数列，但非无穷大量，证明：存在 \(\left\{x_{n}\right\}\) 的两个子列，其中一个子列是无穷大量，另一个是收敛子列．

\subsection{题目2}
2．（15 分）应用 Bolzano－Weierstrass 定理证明闭区间上连续函数的有界性定理．

\subsection{题目3}
3．（10 分）设 \(f_{+}^{\prime}\left(x_{0}\right)>0, f_{-}^{\prime}\left(x_{0}\right)<0\) ，证明：\(x_{0}\) 是 \(f(x)\) 的极小值点．

\subsection{题目4}
4．（10 分）设函数 \(f(x)= \begin{cases}x^{2}, & x \text { 为有理数；} \\ 0, & x \text { 为无理数．证明：} f(x) \text { 只在 } x=0 \text { 处可导．}\end{cases}\)

\subsection{题目5}
5．（10 分）讨论反常积分 \(\int_{0}^{+\infty} \frac{\ln (1+x)}{x^{n}} \mathrm{~d} x\) 的敛散性，其中 \(n\) 为实数．

\subsection{题目6}
6．（15 分）讨论下列函数项级数在所指定区间的一致收敛性：\(\sum_{n=1}^{\infty} 2^{n} \sin \frac{1}{3^{n} x}, x \in(0,+\infty)\) ．

\subsection{题目7}
7．（15 分）判断序列 \(f_{n}(x)=\frac{n x}{1+n^{2} x^{2}}(n=1,2, \cdots)\) 在 \([0,1]\) 上的一致收敛性．

\subsection{题目8}
8．（15 分）设函数 \(f(x, y)\) 在 \(x \in[a, A], y \in[b, B]\) 上连续，函数列 \(\left\{\varphi_{n}(x)\right\}\) 在 \([a, A]\) 上一致收敛，且有 \(b \leq \varphi_{n}(x) \leq B\) ，证明：\(F_{n}(x)=f\left(x, \varphi_{n}(x)\right)\) 在 \([a, A]\) 上一致收敛。

\subsection{题目9}
9．（10 分）设二元函数 \(f(x, y)=\varphi(|x y|)\) ，其中 \(\varphi(0)=0\) ，且在 \(u=0\) 附近满足 \(|\varphi(u)| \leq u^{2}\) ，证明： \(f(x, y)\) 在原点 \((0,0)\) 处可微．

\subsection{题目10}
10．（10 分）计算 \(\iint_{D}|\sin (x-y)| \mathrm{d} x \mathrm{~d} y\) ，其中 \(D=\{(x, y) \mid 0 \leq x \leq y \leq 2 \pi\}\) ．

\subsection{题目11}
11．（15 分）设 \(f(x)\) 在 \([0,+\infty)\) 上连续，且 \(\int_{0}^{+\infty} f(x) \mathrm{d} x\) 收敛，证明：存在 \(\left\{x_{n}\right\} \subset[0,+\infty)\) ，满足\(\lim _{n \rightarrow \infty} x_{n}=+\infty\) ，且 \(\lim _{n \rightarrow \infty} f\left(x_{n}\right)=0\) ．\\
（另一个版本是：设 \(f(x)\) 在 \([0,+\infty)\) 上可导，且 \(\int_{0}^{+\infty} f(x) \mathrm{d} x\) 收敛，证明：存在 \(\left\{x_{n}\right\} \subset[0,+\infty)\) ，满足\(\lim _{n \rightarrow \infty} x_{n}=+\infty\) ，且 \(\lim _{n \rightarrow \infty} f^{\prime}\left(x_{n}\right)=0\) ）

\subsection{题目12}
12．（15 分）计算曲面积分

\[
\begin{gathered}
\iint_{S}(x-y) \mathrm{d} x \mathrm{~d} y+x(y-z) \mathrm{d} y \mathrm{~d} z . \\
\left(\text { 另一个版本为 } \iint_{S} x y^{2} \mathrm{~d} x \mathrm{~d} y+x(y-z) \mathrm{d} y \mathrm{~d} z\right)
\end{gathered}
\]

其中 \(S\) 是圆柱面 \(x^{2}+y^{2}=1\) 被 \(z=0, z=3\) 所微部分的外侧。
\end{document}
